%------------------------------------------------------------
%------------    Estrutura do texto   -----------------------         

% Pacotes Básicos:
%\usepackage{lmodern}			    % Usa a fonte Latin Modern			

\usepackage[T1]{fontenc}		  % Selecao de codigos de fonte.
\usepackage[utf8]{inputenc}		% Codificacao do documento (conversão automática dos acentos)
\usepackage{mathptmx,helvet,courier}
%\usepackage{pslatex}                % Usa a fonte Times New Roman
\usepackage{lastpage}			    % Usado pela Ficha catalográfica
\usepackage[table]{xcolor}
\usepackage{color}				    % Controle das cores
\usepackage{graphicx}			    % Inclusão de gráficos
\usepackage{microtype} 		  	% para melhorias de justificação

\usepackage[normalem]{ulem}

\usepackage{quoting} % usar ambiente quoting (citações longas)

\usepackage{tablefootnote} % para colocar footnotes em tabelas e figuras
\usepackage{caption}
% Pacotes Extras:
\usepackage[final]{pdfpages}
\usepackage{amsmath,amsthm}   %Símbolos Matemáticos
%\usepackage{indentfirst} % Indenta primeiro parágrafo 
\usepackage[portuguese, ruled, linesnumbered,commentsnumbered, algo2e, vlined, lined, boxed, algochapter]{algorithm2e} % Algoritmos 
\usepackage{hyperref}
\usepackage[brazilian,hyperpageref]{backref}	 % Paginas com as citações na bibliograficas
% Configurações do pacote backref
% Usado sem a opção hyperpageref de backref
\renewcommand{\backrefpagesname}{Citado na(s) página(s):~}
% Texto padrão antes do número das páginas
\renewcommand{\backref}{}
% Define os textos da citação
\renewcommand*{\backrefalt}[4]{
\ifcase #1 %
Nenhuma citação no texto.%
\or
Citado na página #2.%
\else
Citado
\fi}%


%\usepackage[num]{abntex2cite}	% Citações padrão ABNT númerico
\usepackage[alf,abnt-emphasize=bf,abnt-full-initials=no,abnt-etal-list=5,abnt-etal-text=it,abnt-repeated-author-omit=yes]{abntex2cite}
%\usepackage[alf,abnt-emphasize=bf,abnt-etal-list=0,abnt-etal-text=emph]{abntex2cite}

% se desejar justificar as referências (pela abnt é alinhado a esquerda)
% \usepackage[alf,abnt-emphasize=bf,bibjustif]{abntex2cite}

\usepackage{setspace}

\usepackage{float} % para ajustar a posição das imagens

\usepackage{tikzsymbols} %para caracteres emoji
\usepackage{stackengine}
\usepackage{scalerel}


\newcommand\dangersign[1][2ex]{%
  \renewcommand\stacktype{L}%
  \scaleto{\stackon[1.3pt]{\color{red}$\triangle$}{\tiny\bfseries !}}{#1}%
}

\usepackage{etoolbox}
\usepackage{lipsum}
\usepackage{amssymb}
\usepackage{amsmath}

%para usar os algoritmos e matematica
\newcommand{\var}{\texttt}

% Defininfo Cores:
\definecolor{blue}{RGB}{25,25,112}
\definecolor{midgray}{gray}{.7}

\makeatletter % informações do PDF
\hypersetup{ % pagebackref=true,
	pdftitle={\@title}, 
	pdfauthor={\@author},
    pdfsubject={\imprimirpreambulo},
	pdfcreator={LaTeX with abnTeX2},
	pdfkeywords={abnt}{latex}{abntex}{abntex2}{trabalho acadêmico}, 
	colorlinks=true,     % false: boxed links; true: colored links
    linkcolor=black,          	% color of internal links
    citecolor=black,        		% color of links to bibliography
    filecolor=magenta,      	% color of file links
    urlcolor=black,
	bookmarksdepth=4 }
\makeatother
 

% -------------------------------------------- 
% Espaçamentos entre linhas e parágrafos 
\setlength{\parindent}{1.3cm} % O tamanho do parágrafo

% Controle do espaçamento entre um parágrafo e outro:
\setlength{\parskip}{0.2cm}  % tente também \onelineskip

% Definição de ambientes matemáticos em português 
\newtheorem{teorema}{Teorema}[chapter]
\newtheorem{axioma}{Axioma}[chapter]
\newtheorem{corolario}{Corolário}[chapter]
\newtheorem{lema}{Lema}[chapter]
\newtheorem{proposicao}{Proposição}[chapter]
\newtheorem{definicao}{Definição}[chapter]
\newtheorem{exemplo}{Exemplo}[chapter]
\newtheorem{observacao}{Observação}[chapter]

% Novos Comandos
\usepackage{tgtermes}

% Evita com que todas as seções fiquem com o estilo do título
\renewcommand{\tocheadstart}{\normalshape}

%%%%%%%%%%%%%%%%%%%%%%%%%%%%%%%%%%%%%%%%%%%%%%%%%%%%%%%%%%
% mostra no máximo subsubsection no sumário
\settocdepth{subsubsection} 

% diminuir espaço entre número e título do capítulo/seção/etc
\addtolength{\cftlastnumwidth}{-1.5em}
\cftsetindents{section}{0em}{\cftlastnumwidth}
\cftsetindents{subsection}{0em}{\cftlastnumwidth}
\cftsetindents{subsubsection}{0em}{\cftlastnumwidth}

% números de seções sem qualquer destaque
% seção primária (capítulo)
\let\oldthechapter\thechapter
\renewcommand{\thechapter}{\textnormal{\oldthechapter}}

% seção secundária (section)
\let\oldthesection\thesection
\renewcommand{\thesection}{\textnormal{\oldthesection}}

% seção terciária (subsection)
\let\oldthesubsection\thesubsection
\renewcommand{\thesubsection}{\textnormal{\oldthesubsection}}

% seção quaternária (subsubsection)
\let\oldthesubsubsection\thesubsubsection
\renewcommand{\thesubsubsection}{\textnormal{\oldthesubsubsection}}

% seção quinária (subsubsubsection)
\let\oldthesubsubsubsection\theparagraph
\renewcommand{\theparagraph}{\textnormal{\oldthesubsubsubsection}}


%%%%%%%%%%%%%%%%%%%%%%%%%%%%%%%%%%%%%%%%%%%%%%%%%%%%%%%%%%
% Estilos das seções

% seção primária (capítulo) em letra maiúscula, em negrito
\renewcommand{\ABNTEXchapterfont}{\rmfamily\bfseries}
\renewcommand{\ABNTEXchapterfontsize}{\normalsize}
\setboolean{ABNTEXupperchapter}{true}

% seção secundária em negrito 
\renewcommand{\ABNTEXsectionfont}{\rmfamily\bfseries}
\renewcommand{\ABNTEXsectionfontsize}{\normalsize}
\setboolean{ABNTEXuppersection}{false}

% seção terciária negrito e itálico
\renewcommand{\ABNTEXsubsectionfont}{\rmfamily\bfseries\itshape}
\renewcommand{\ABNTEXsubsectionfontsize}{\normalsize}
\setboolean{ABNTEXuppersubsection}{false}

% seção quaternária normal
\renewcommand{\ABNTEXsubsubsectionfont}{\rmfamily}
\renewcommand{\ABNTEXsubsubsectionfontsize}{\normalsize}
\setboolean{ABNTEXuppersubsubsection}{false}

% seção quinária em itálico
\renewcommand{\ABNTEXsubsubsubsectionfont}{\rmfamily\itshape}
\renewcommand{\ABNTEXsubsubsubsectionfontsize}{\normalsize}
\setboolean{ABNTEXuppersubsubsubsection}{false}

%%%%%%%%%%%%%%%%%%%%%%%%%%%%%%%%%%%%%%%%%%%%%%%%%%%%%%%
% Configuração das fontes no sumário

\renewcommand{\cftchapterfont}{\ABNTEXchapterfont}
\renewcommand{\cftchapterpagefont}{\normalsize}

\renewcommand{\cftsectionfont}{\ABNTEXsectionfont}
\renewcommand{\cftsectionpagefont}{\normalsize}

\renewcommand{\cftsubsectionfont}{\ABNTEXsubsectionfont}
\renewcommand{\cftsubsectionpagefont}{\normalsize}
  
\renewcommand{\cftsubsubsectionfont}{\ABNTEXsubsubsectionfont}
\renewcommand{\cftsubsubsectionpagefont}{\normalsize}

\renewcommand{\cftparagraphfont}{\ABNTEXsubsubsubsectionfont}
\renewcommand{\cftparagraphpagefont}{\normalsize}

% legenda de figuras com fonte reduzida
\let\oldlegend\legend
\renewcommand{\legend}[1]{\oldlegend{\ABNTEXfontereduzida{#1}}}

% Variáveis adicionais
\providecommand{\imprimirautorcite}{}
\newcommand{\autorcite}[1]{\renewcommand{\imprimirautorcite}{#1}} 
\providecommand{\imprimirsigla}{}
\newcommand{\sigla}[1]{\renewcommand{\imprimirsigla}{#1}}
\providecommand{\imprimiruf}{}
\newcommand{\uf}[1]{\renewcommand{\imprimiruf}{#1}}
\providecommand{\imprimircurso}{}
\newcommand{\curso}[1]{\renewcommand{\imprimircurso}{#1}}
\providecommand{\imprimirinstituto}{}
\newcommand{\instituto}[1]{\renewcommand{\imprimirinstituto}{#1}}
\providecommand{\imprimirdepartamento}{}
\newcommand{\departamento}[1]{\renewcommand{\imprimirdepartamento}{#1}}
\providecommand{\imprimirano}{}
\newcommand{\ano}[1]{\renewcommand{\imprimirano}{#1}}
\providecommand{\imprimirdia}{}
\newcommand{\dia}[1]{\renewcommand{\imprimirdia}{#1}}
\providecommand{\imprimirmes}{}
\newcommand{\mes}[1]{\renewcommand{\imprimirmes}{#1}}
\providecommand{\imprimirgrau}{}
\newcommand{\grau}[1]{\renewcommand{\imprimirgrau}{#1}}
\providecommand{\imprimirexaminadorum}{}
\newcommand{\examinadorum}[1]{
    \renewcommand{\imprimirexaminadorum}{#1}}
\providecommand{\imprimirexaminadordois}{}
\newcommand{\examinadordois}[1]{
    \renewcommand{\imprimirexaminadordois}{#1}}
\providecommand{\imprimirexaminadortres}{}
\newcommand{\examinadortres}[1]{
    \renewcommand{\imprimirexaminadortres}{#1}}
\providecommand{\imprimirexaminadorquatro}{}
\newcommand{\examinadorquatro}[1]{
    \renewcommand{\imprimirexaminadorquatro}{#1}}
\providecommand{\imprimirttorientador}{}
\newcommand{\ttorientador}[1]{
    \renewcommand{\imprimirttorientador}{#1}} 
\providecommand{\imprimirttcoorientador}{}
\newcommand{\ttcoorientador}[1]{
    \renewcommand{\imprimirttcoorientador}{#1}}
\providecommand{\imprimirttexaminadorum}{}
\newcommand{\ttexaminadorum}[1]{
    \renewcommand{\imprimirttexaminadorum}{#1}}
\providecommand{\imprimirttexaminadordois}{}
\newcommand{\ttexaminadordois}[1]{\renewcommand{
        \imprimirttexaminadordois}{#1}}
\providecommand{\imprimirttexaminadortres}{}
\newcommand{\ttexaminadortres}[1]{
    \renewcommand{\imprimirttexaminadortres}{#1}}
\providecommand{\imprimirttexaminadorquatro}{}
\newcommand{\ttexaminadorquatro}[1]{
    \renewcommand{\imprimirttexaminadorquatro}{#1}}

% Cria o comando \subtitulo. A norma define que o TITULO deve ser em caixa alta
% negrito, mas o subtitulo deve ser em caixa baixa.
\providecommand{\imprimirsubtitulo}{}
\newcommand{\subtitulo}[1]{\renewcommand{\imprimirsubtitulo}{#1}}

%----------------------------------------------------
\renewcommand{\imprimircapa}{  % Capa 
\begin{capa}
\center
\MakeTextUppercase{\imprimirinstituicao}\\ \MakeTextUppercase{\imprimirinstituto} \\
%\imprimirdepartamento \\
\MakeTextUppercase{\imprimircurso} \\              
\vspace{5cm}
\imprimirautor\\
\vspace{3.5cm}
\MakeTextUppercase{\textbf{\imprimirtitulo}}  \\
\MakeTextLowercase{\textbf{\imprimirsubtitulo}}
\vspace{2cm}
\vfill
\imprimirlocal~-~\imprimiruf \\ 
\imprimirano 
\end{capa}   } % Capa

%----------------------------------------------------
\renewcommand{\imprimirfolhaderosto}{% folha de rosto
    \begin{center}
    {{\MakeTextUppercase \imprimirautor}}  \\
		\vfill
		{\textbf{\MakeTextUppercase{\imprimirtitulo}}\\
        \MakeTextLowercase{\textbf{\imprimirsubtitulo}}}
    \end{center}
    \vfill 
    \begin{flushright} 
    \parbox{0.48\linewidth}{
	\imprimirtipotrabalho~ apresentado ao curso de Bacharelado em \imprimircurso~ da \imprimirinstituicao~ - \imprimirinstituto~ para a obtenção do título de \imprimirgrau. \\
	\vfill
	{\imprimirorientadorRotulo}~\imprimirorientador \\
	\vfill 
	{\imprimircoorientadorRotulo}~\imprimircoorientador}
   \end{flushright} 
	 \vfill
   \begin{center}
   {\large{\imprimirlocal~- \imprimiruf \\ \imprimirano}}
   \end{center} }  % folha de rosto

%----------------------------------------------------


\renewcommand{\listfigurename}{LISTA DE ILUSTRAÇÕES}
\pdfbookmark[0]{\listfigurename}{lof}

%================================================================================
% Pacotes de citações
%================================================================================

% ---
\makeatletter
\newcommand\thefontsize{Fonte tamanho: \f@size pt}
\makeatother



% para não deixar uma figura sozinha no centro de uma página,
% mas no topo da página.
\makeatletter
\setlength{\@fptop}{0pt}
\makeatother

\DeclareMathOperator*{\argmin}{arg\,min}

\let\theoldtextual\textual
\renewcommand{\textual}{
\theoldtextual
\pagenumbering{arabic} 
\setcounter{page}{\thesheetsequence-1}
}

\newcommand{\listafiguras}{
\listoffigures*   % Cria a Lista de Figuras
\cleardoublepage
}

% ---
\newcommand{\listatabelas}{
\pdfbookmark[0]{\listtablename}{lot}
\listoftables*  % Cria a lista de Tabelas
\cleardoublepage
}

\renewcommand{\listalgorithmcfname}{Lista de algoritmos}
\newcommand{\listaalgoritmos}{
\pdfbookmark[0]{\listalgorithmcfname}{lof}
\listofalgorithmes   % Cria a lista de Tabelas
\cleardoublepage
}

\newcommand{\sumario}{
\pdfbookmark[0]{\contentsname}{toc}
\tableofcontents*
\cleardoublepage
}

%--